\section{Spin slots, loops, and pins}\label{sec:loops}

We first express the finite-volume Gibbs state by exchange paths.  The
representation will identify the entire spectral distribution of total
magnetization, which is needed when a small field is introduced at the end
of the proof.  We then define the pin laws used in the intervening
probabilistic argument and establish their two elementary measure bounds.
Powers related the isotropic Heisenberg model to random walks on
permutations~\cite{Powers1976}.  T\'oth gave the random-stirring
representation for spin $1/2$~\cite{Toth1993}.  Nachtergaele developed
general-spin symmetric-slot representations~\cite[Section~2.2,
equations~(2.23)--(2.25), and Theorem~2.4]{Nachtergaele1994Quasistates};
see also his account of geometric representations~\cite{Nachtergaele1994Stochastic}.
For the endpoint-permutation realization and the magnetization
transform, see also Bj\"ornberg, Fr\"ohlich and
Ueltschi~\cite[Section~3.2, Proposition~3.2 and equation~(3.19)]{BjornbergFrohlichUeltschi2020}.
Their representation is presented for the complete graph with a different
coupling normalization.  We derive the finite-graph algebra here for
our nearest-neighbor Hamiltonian, including its rates and full spectral
law, before introducing the pin weights needed in the proof.

Fix a finite cube $\Lambda\subset\Z^d$, $S\in\{\tfrac12,1,\tfrac32,\ldots\}$,
and $0<\beta<\infty$.  Set $\ell=2S$ and introduce the \emph{slot set}
\[
 V_\Lambda=\Lambda\times\{1,\ldots,\ell\}.
\]
A slot $i=(x,a)$ is located at the site $x$.  Two slots form a
\emph{slot edge} if their sites are nearest neighbors; all $\ell^2$ pairs
above each lattice edge are included.  Let $\mathcal E_\Lambda$ be this
set of unordered slot edges and put $q_\Lambda=|\mathcal E_\Lambda|$.
On
\[
 \widehat{\mathcal H}_\Lambda
   =\bigotimes_{i\in V_\Lambda}\C^2
\]
write $s_i^a$ for the spin-$\tfrac12$ matrix in direction $a$ acting on
slot $i$.  For a permutation $\pi$ of the slots, $\mathcal U_\pi$ denotes
the tensor permutation sending the factor originally at $i$ to
$\pi(i)$; in particular
$\mathcal U_\pi\mathcal U_\rho=\mathcal U_{\pi\circ\rho}$.

\begin{lemma}[Symmetric-slot realization]\label{lem:slot-realization}
At each site $x$, the operators
$\widehat S_x^a=\sum_{j=1}^{\ell}s_{(x,j)}^a$ restrict on
$\operatorname{Sym}^{\ell}(\C^2)$ to the spin-$S$ matrices defined in
Section~\ref{sec:model}.  The orthogonal projection onto this subspace is
\[
 \mathcal S_x=\frac1{\ell!}
      \sum_{\kappa\in\mathfrak S_\ell}\mathcal U_{\kappa,x}.
\]
With $\mathcal S_\Lambda=\prod_{x\in\Lambda}\mathcal S_x$, define
\[
 \widehat H_\Lambda
 =-\sum_{\{x,y\}\subset\Lambda:\,|x-y|_1=1}
          \sum_{a=x,y,z}\widehat S_x^a\widehat S_y^a,
 \qquad
 \widehat M_\Lambda=\sum_{i\in V_\Lambda}s_i^z.
\]
Both operators commute with every within-site permutation and with
$\mathcal S_\Lambda$.  Their restrictions to
$\operatorname{ran}\mathcal S_\Lambda$ are unitarily equivalent to
$H_\Lambda$ and $M_\Lambda=\sum_{x\in\Lambda}S_x^z$, respectively.
Moreover,
\begin{equation}\label{eq:slot-hamiltonian}
 \widehat H_\Lambda
   =-\frac12\sum_{\{i,j\}\in\mathcal E_\Lambda}
                   \mathcal U_{(ij)}
       +\frac{q_\Lambda}{4}I.
\end{equation}
\end{lemma}

\begin{proof}
Let $\psi_k$ be the normalized sum of the $\binom\ell k$ elementary
tensor vectors having exactly $k$ up spins, for $0\le k\le\ell$.
These form an orthonormal basis of the symmetric subspace, and
\[
 \widehat S^z\psi_k=(k-\ell/2)\psi_k.
\]
For $0\le k<\ell$, each elementary vector with $k+1$ up spins occurs
$k+1$ times in $\sum_j s_j^+\psi_k$.  Consequently
\[
 \widehat S^+\psi_k
  =(k+1)\sqrt{\frac{\binom\ell{k+1}}{\binom\ell k}}\,\psi_{k+1}
  =\sqrt{(\ell-k)(k+1)}\,\psi_{k+1}.
\]
Also $\widehat S^+\psi_\ell=0$.  Under $m=k-S$, the squared coefficient is
$S(S+1)-m(m+1)$.  The lowering, $x$, and $y$ matrices follow by
adjunction and the definitions.  Thus $e_m\mapsto\psi_{S+m}$ gives the
claimed unitary identification.

The average over the finite permutation group is self-adjoint and
idempotent, and its range is exactly the invariant subspace, proving the
projection formula.  Conjugation by a within-site permutation merely
reorders the summands of each $\widehat S_x^a$.  It therefore commutes
with $\widehat H_\Lambda$ and $\widehat M_\Lambda$.  Equivalently, in
the slot-edge expression, conjugation permutes the complete collection
of $\ell^2$ slot pairs above every incident lattice edge.  This last
observation is essential: an individual inter-site transposition need
not commute with $\mathcal S_x$.

Finally, on two spin-$\tfrac12$ factors, direct calculation in the
up/down basis gives
\[
 \sum_{a=x,y,z}s_i^as_j^a
       =\frac12\mathcal U_{(ij)}-\frac14 I.
\]
Summing this identity over the slot edges proves
\eqref{eq:slot-hamiltonian} and the asserted restrictions.
\end{proof}

The \emph{base exchange law} $\mathbb P_{\mathrm b}$ is defined as
follows.  Independently on each edge of $\mathcal E_\Lambda$, put a
Poisson process of rate $1/2$ on the time interval $[0,\beta)$.
At a mark, exchange the two slot lines.  At the end of the interval,
apply an independent uniform permutation of the $\ell$ slots at each
site, and then identify time $\beta$ with time $0$.
An \emph{exchange picture} $\omega$ comprises these marks and endpoint
permutations.  Following its lines forward through repeated copies of
the same picture gives directed loops, which we call its \emph{cycles}.
They correspond to the cycles of its one-period slot permutation
$\pi_\omega$ and can wind through several periods.  Let
$c(\omega)$ be their number.  In particular,
$c(\omega)\le |V_\Lambda|$.

Slot-time points belong to
$V_\Lambda\times(\R/\beta\Z)$.  At an exchange we use the slot just
after the exchange; time $0$ is understood just after the within-site
endpoint permutations and the wrapping of time.  Thus every cycle has
a well-defined set of slots at the time-zero cut.  Write $a_C$ for the
number of these slots on a cycle $C$; then $\sum_C a_C=|V_\Lambda|$.

\begin{proposition}[Exact spin-to-loop representation]\label{prop:spin-loops}
For every finite cube $\Lambda$, spin $S$, and $0<\beta<\infty$ as
above, the partition function is
\begin{equation}\label{eq:loop-partition-function}
 Z_\Lambda(\beta):=\Tr_{\mathcal H_\Lambda}e^{-\beta H_\Lambda}
       =e^{\beta q_\Lambda/4}\,
                     \mathbb E_{\mathrm b}2^{c(\omega)}.
\end{equation}
More generally, for every function $\Phi$ on the finite spectrum of
$M_\Lambda$,
\begin{equation}\label{eq:spectral-loop-law}
 \frac{\Tr_{\mathcal H_\Lambda}
          \bigl(e^{-\beta H_\Lambda}\Phi(M_\Lambda)\bigr)}
      {Z_\Lambda(\beta)}
 =\frac{\displaystyle\mathbb E_{\mathrm b}
          \sum_{\varepsilon\in\{-1,1\}^{\{\text{cycles of }\omega\}}}
          \Phi\!\left(\frac12\sum_C\varepsilon_C a_C\right)}
        {\displaystyle\mathbb E_{\mathrm b}2^{c(\omega)}}.
\end{equation}
Thus the spectral law of $M_\Lambda$ is obtained by weighting the base
picture by $2^{c(\omega)}$, then coloring each cycle independently
up or down with equal probabilities, and summing the corresponding
$+1/2$ or $-1/2$ over its slots at time $0$.
\end{proposition}

\begin{proof}
Let $\varphi_t$ be the permutation generated by the inter-site marks
up to time $t$, without the endpoint permutations.  Independent
increments, or the Poisson expansion in time-ordered marks, give
\[
 \mathbb E_{\mathrm b}\mathcal U_{\varphi_\beta}
   =\exp\!\left\{
       \frac\beta2\sum_{\{i,j\}\in\mathcal E_\Lambda}
                    (\mathcal U_{(ij)}-I)\right\}
   =e^{-\beta q_\Lambda/4}e^{-\beta\widehat H_\Lambda}.
\]
Indeed, over an interval of length $\dd t$, the expectation is
$I+\frac12\sum_{\{i,j\}}(\mathcal U_{(ij)}-I)\dd t+o(\dd t)$;
the resulting constant-coefficient matrix differential equation gives
the displayed exponential.  Averaging the independent endpoint
permutations now gives
\begin{equation}\label{eq:loop-operator-average}
 \mathbb E_{\mathrm b}\mathcal U_{\pi_\omega}
    =e^{-\beta q_\Lambda/4}
                   \mathcal S_\Lambda e^{-\beta\widehat H_\Lambda}.
\end{equation}
Lemma~\ref{lem:slot-realization} identifies the trace on the range of
$\mathcal S_\Lambda$ with the physical trace.  Moreover,
$\widehat M_\Lambda$ commutes with every slot permutation, since it
is the sum of the same one-slot observable over all slots.  Therefore
\eqref{eq:loop-operator-average} yields
\begin{equation}\label{eq:loop-spectral-trace}
 \Tr_{\mathcal H_\Lambda}
       \bigl(e^{-\beta H_\Lambda}\Phi(M_\Lambda)\bigr)
 =e^{\beta q_\Lambda/4}\,
   \mathbb E_{\mathrm b}\Tr_{\widehat{\mathcal H}_\Lambda}
       \bigl(\mathcal U_{\pi_\omega}\Phi(\widehat M_\Lambda)\bigr).
\end{equation}

To calculate the last trace, index the elementary up/down tensor
basis by $\eta\in\{-1,1\}^{V_\Lambda}$.  The diagonal entry of
$\mathcal U_{\pi_\omega}$ at $\eta$ is one exactly when $\eta$ is
constant on every cycle of $\pi_\omega$, and is zero otherwise.
For such a basis vector the eigenvalue of $\widehat M_\Lambda$ is
$\frac12\sum_C\varepsilon_Ca_C$.  Hence
\[
 \Tr_{\widehat{\mathcal H}_\Lambda}
       \bigl(\mathcal U_{\pi_\omega}\Phi(\widehat M_\Lambda)\bigr)
   =\sum_{\varepsilon\in\{-1,1\}^{\{\text{cycles of }\omega\}}}
          \Phi\!\left(\frac12\sum_C\varepsilon_Ca_C\right).
\]
Taking $\Phi=1$ proves \eqref{eq:loop-partition-function}; dividing
\eqref{eq:loop-spectral-trace} by it proves
\eqref{eq:spectral-loop-law}.  In particular, taking $\Phi$ to be
the indicator of any spectral subset proves the claim about the
spectral distribution, not only its moments.
\end{proof}

\subsection{Pin laws and the boundary limit}

The factor two in the representation records the two colors of each
cycle.  A pin fixes the color of every cycle it meets.  This motivates
a family of measures in which only the cycles avoiding the pins
retain their factor two.

For a deterministic pin set
$D\subset V_\Lambda\times(\R/\beta\Z)$, let
\[
 n_D(\omega)=\#\{C:C\text{ is a cycle of }\omega,\ C\cap D=\varnothing\},
 \qquad w_D(\omega)=2^{n_D(\omega)},
 \qquad Z_D=\mathbb E_{\mathrm b}w_D.
\]
We use finite pin sets and their unions with the continuous boundary
set defined below.  Their cycle-incidence events are measurable.
The \emph{pin law} and its expectation are
\begin{equation}\label{eq:pin-law}
 \mathbb P_D(\dd\omega)
      =\frac{w_D(\omega)}{Z_D}\mathbb P_{\mathrm b}(\dd\omega),
 \qquad \mathbb E_D[\cdot].
\end{equation}
The set of slot-time points on cycles missing $D$ is denoted by
\[
 \mathcal B_D(\omega)=\bigcup_{C:\,C\cap D=\varnothing}C.
\]
In fixed volume, $1\le w_D\le 2^{|V_\Lambda|}$, so $Z_D$ is finite
and strictly positive.  The empty-pin law $\mathbb P_\varnothing$
is exactly the picture marginal in Proposition~\ref{prop:spin-loops}.

The \emph{inner boundary} and the continuous boundary pin set are
\begin{equation}\label{eq:continuous-boundary-pins}
 \partial\Lambda
    =\{x\in\Lambda:\exists y\notin\Lambda,\ |x-y|_1=1\},
 \qquad
 D_0=\partial\Lambda\times\{1,\ldots,\ell\}
                       \times(\R/\beta\Z).
\end{equation}
Every slot above a boundary site is therefore pinned at every time.
The following fixed-volume approximation allows finite-layer
arguments to be applied with this boundary condition.

\begin{lemma}[Dense boundary pins]\label{lem:dense-pins}
Fix a finite cube $\Lambda$, $0<\beta<\infty$, and a finite
deterministic pin set $F$.  Let $(\mathcal C_n)_{n\ge1}$ be increasing
finite sets of times, each containing $0$, whose union is dense in
$\R/\beta\Z$, and put
\[
 D_{0,n}=\partial\Lambda\times\{1,\ldots,\ell\}\times\mathcal C_n,
 \qquad D_n=F\cup D_{0,n},\qquad D=F\cup D_0.
\]
For $\mathbb P_{\mathrm b}$-almost every picture there is a finite
$n_*(\omega)$ such that, for all $n\ge n_*(\omega)$,
\[
 w_{D_n}(\omega)=w_D(\omega),\qquad
 \mathcal B_{D_n}(\omega)=\mathcal B_D(\omega).
\]
The laws $\mathbb P_{D_n}$ converge to $\mathbb P_D$ in total
variation.  More generally, if $\Phi_n\to\Phi$ almost surely under
$\mathbb P_{\mathrm b}$ and $\sup_n|\Phi_n|\le A<\infty$, then
\begin{equation}\label{eq:dense-pin-expectations}
             \mathbb E_{D_n}\Phi_n\longrightarrow\mathbb E_D\Phi.
\end{equation}
\end{lemma}

\begin{proof}
In one period the number of Poisson marks is finite almost surely.
Their times are distinct and avoid the countable collection of
deterministic cut times and pin times, almost surely.  A line that
visits a boundary site occupies that site on a time interval of
positive length: inter-site marks do not coincide, and endpoint
permutations preserve the site.  This remains true for a visit
crossing the time seam.  A dense collection of cuts detects every
such visit, with every slot included at each boundary cut.

There are only finitely many cycles.  For a fixed picture, choose
one detecting cut for each cycle that meets $D_0$ and take the
largest of their finite indices.  After this index every cycle has
the same incidence with $D_n$ as with $D$.  This proves both eventual
equalities.  Also $w_{D_n}\downarrow w_D$ and
$w_{D_n}\le 2^{|V_\Lambda|}$.  Dominated convergence gives
$Z_{D_n}\to Z_D$ and
\[
 \mathbb E_{\mathrm b}
       \left|\frac{w_{D_n}}{Z_{D_n}}-\frac{w_D}{Z_D}\right|
       \longrightarrow0,
\]
which proves total-variation convergence.  Applying dominated
convergence to $w_{D_n}\Phi_n$ proves
\eqref{eq:dense-pin-expectations}.  The dominating constant here
depends on the fixed volume; no uniform-in-volume bound on the
total loop weight is asserted.
\end{proof}

\subsection{Factorial domination for exchange marks}

The pin weights alter the distribution of the Poisson marks.  The
bound needed later is on ordered lists of distinct marks; it does
not require independence under a pin law.

Let
$X_\Lambda=\mathcal E_\Lambda\times[0,\beta)$ and let
$\mu$ be the measure satisfying
$\mu(\{e\}\times\dd t)=\tfrac12\dd t$.
Write $\omega_{\mathrm m}$ for the counting measure of inter-site
marks in the picture.  For $n\ge1$, the notation
$(\omega_{\mathrm m})^n_{\ne}$ means ordered $n$-tuples of distinct
marks, even when some of their edge labels agree.

\begin{lemma}[Factorial-measure domination]\label{lem:factorial}
For every finite pin set $D$, or such a set united with $D_0$, every
integer $n\ge1$, and every nonnegative measurable function
$F:X_\Lambda^n\to[0,\infty]$,
\begin{equation}\label{eq:factorial-domination}
 \mathbb E_D\sum_{(u_1,\ldots,u_n)\in(\omega_{\mathrm m})^n_{\ne}}
                 F(u_1,\ldots,u_n)
 \le 2^n\int_{X_\Lambda^n}F(u_1,\ldots,u_n)
                          \prod_{j=1}^n\mu(\dd u_j).
\end{equation}
In particular, the estimate holds with the marks restricted to any
specified time interval and with any specified chronological order.
\end{lemma}

\begin{proof}
Inserting one inter-site transposition into a picture cuts two
directed strands and reconnects them with their outgoing portions
exchanged.  If they lie on different cycles, those cycles join; if
they lie on one cycle, that cycle splits into two.  The union of
their slot-time strands, and hence all pin incidences on that
union, is retained.  On joining, the number of cycles missing $D$
cannot increase.  On splitting, it increases by at most one: a
cycle missing $D$ gives two such cycles, while a cycle meeting $D$
has at least one descendant still meeting $D$.  Thus, away from
coincident mark times,
\begin{equation}\label{eq:pin-insertion-ratio}
 w_D\!\left(\omega+\sum_{j=1}^n\delta_{u_j}\right)
       \le 2^n w_D(\omega).
\end{equation}
For finite pins coincidences with pin times have product-intensity
measure zero; including $D_0$ does not change the surgery argument.
The endpoint permutations are held fixed throughout this insertion.

The identity used to integrate \eqref{eq:pin-insertion-ratio} is the
iterated Mecke formula~\cite{Mecke1967}; its factorial form is stated in
\cite[Theorem~4.4]{LastPenrose2017}.  We derive it here under the base
Poisson law, with the pin weight retained inside the integrand.  Write $\lambda=\mu(X_\Lambda)$.
For a fixed endpoint permutation, expand the Poisson measure as
\[
 e^{-\lambda}\sum_{k=0}^{\infty}\frac1{k!}
                      \prod_{j=1}^k\mu(\dd u_j).
\]
In the weighted sum over ordered distinct $n$-tuples, the $k$-mark
term has $k!/(k-n)!$ choices of the distinguished marks.  Relabel
them $u_1,\ldots,u_n$ and sum the remaining $k-n$ marks.  Averaging
also over the independent endpoint permutations gives the exact
identity
\begin{align*}
 &\mathbb E_{\mathrm b}\left[
       w_D(\omega)
       \sum_{(u_1,\ldots,u_n)\in(\omega_{\mathrm m})^n_{\ne}}
                    F(u_1,\ldots,u_n)\right]\\
 &\hspace{8mm}=
   \int_{X_\Lambda^n}F(u_1,\ldots,u_n)
       \mathbb E_{\mathrm b}\!\left[
          w_D\!\left(\omega+\sum_{j=1}^n\delta_{u_j}\right)\right]
             \prod_{j=1}^n\mu(\dd u_j).
\end{align*}
All terms are nonnegative, so this manipulation is justified also
when an integral is infinite.  By \eqref{eq:pin-insertion-ratio},
the expectation inside the integral is at most $2^n Z_D$.
Division by $Z_D$ proves \eqref{eq:factorial-domination}.
\end{proof}

The estimate is unconditional under $\mathbb P_D$.  If an independent
fresh Poisson family is also sampled, mixed lists of distinct actual
and fresh marks satisfy the corresponding product bound: an actual
mark contributes $2\mu$, and a fresh mark contributes $\mu$.
Indeed, fix which entries in the list come from each family, use
independence to factor the two factorial measures, and apply
Lemma~\ref{lem:factorial} to the actual-mark factor.  This observation
will permit chronological path estimates after averaging over an
exchange picture and fresh randomness; it makes no conditional
factorial-domination assertion at a fixed path exposure.

We now have the exact spectral representation and a family of pin
laws with controlled mark counts.  The next section estimates the
probability that many prescribed slot-time points simultaneously
belong to cycles avoiding the existing pins.
